%%
%% Author: shoupu_wan
%% 7/14/18
%%

\chapter{Flatten Nested Loops} \label{ch:loop-thinning}

One of the common mistake that an inexperienced developer tend to make is using
nested loop in situations where it could have been avoided.
Although permitted by most programming languages, nested loops almost always bring about
hindrance to code readability and maintainability.
%More often than not, they are mere nuisances when one tries to understand the semantics
%or analyze the runtime complexity.
Occasionally, they may be the cause of degenerate runtime for suboptimal
implementations---\textit{e.g.}, turning a $O(N)$ algorithm into one of $O(N^2)$. %TODO example?


The difficulty of adopting a new programming habit in this regard is
awareness---when compelled to use a nested loop, developers need to take a step back
and think over the options instead of jump into conviction.
When it comes to algorithmic implementation using nested loops, the rule of thumb is ``Be cautious'.
More often than not, nested loops are not necessary and can be compressed into single-layer,
dynamically paced loops with the help of auxiliary data structures, cascade conditionals, and so on.

This chapter is a spin-off of the ~\autoref{ch:loops} so that
we may dedicate this chapter to the study of simplifying nested loops.


\section{Empirical Loop-thinning Technique}\label{sec:avoidable-scenarios-nested-loop}


%First working version is seldom the best solution.
%Let us ask this question: ``Under what circumstances can nested loops be replaced by thinner
%loops?''
Many situations that nested loops appear so compelling or even inevitable are not so
in retrospection.
Speaking empirically, nested loops may possibly be avoided in these scenarios---algorithmic
implementations where the data is moved down more than one data structures,
like in the settings of an assembly line.


The only scenario where I found myself needing nested loops is when working with $2$D or higher
dimensional data structures like list of lists, matrices, \textit{etc.}, where use of index
tuples such as \code{i,j} seemingly necessitates nested loops.
Even so, there are plenty of opportunities for one to refactor nested loops into more meaningful
functions with meaningful names for better code readability and reusability,
each containing a single loop.
Some may argue that function invocation overhead may hurt performance especially inside a loop.
My answer to that is that if you really care about performance you would have vectorized
your code to use built-in matrix operations instead of nested loops.
%But whichever way you go, it involves function calls.


One particular case that almost comes with guaranteed success is nested loop where on the
outside is an end-to-end, evenly paced loop while on the inside is a loop that does partial
lookup either on an auxiliary data structure or from a `slice' of the same data
structure (for example, an array slice from the starting point to the current index).
When an auxiliary data structure is used, it is used to provide cumulative, partial
information in addition to the main data structure.


In this situation, one or more of the following measures may be taken to eliminate
the nested, inner loop
\begin{itemize}
    \item Use of auxiliary data structures such as a queue or stack to facilitate with
    dynamic paced iteration;
    \item Use of additional conditional construct or additional test to existing ones;
    \item Make loop pacing to be more dynamic;
    \item Other changes necessary for boundary condition handling.
\end{itemize}


The bare bone changes for flattening nested loops is simple---identify the conditions for the
inner loop, convert the inner loop construct to a conditional or piggyback to an existing
conditional (and make cascade conditionals, see ~\autoref{ch:conditional-cascade}),
relocate the state-updating statement(s) of the outer loop to a conditional block
(possibly an \code{else}-block in the newly formed conditional statement).


Other changes may be performed after these in conformation to conventions and syntactic aesthetics.
Most importantly, following the removal of nested loops, other optimizations may be unlocked to
further the code quality to another level.
Like a critical piece is added in a candy crush game which causes chains of avalanches.
That is the true value of of code optimization.


Put together, we refer to the coding strategies to coalesce a nested loop into a
single-layer one as `loop thinning' technique.
In what follows, we are going to demonstrate application of `loop thinning' technique through
a few examples adapted from real-world programming projects.


Because most of the methods we present in this chapter are empirical, I cannot claim
universal applicability for many of them.
Instead, the techniques may be applied on a case by case basis.


\section{Case Study: I}\label{sec:nested-loop-histogram}


Continuing with the example ~\autoref{lst:histogram-oneMoreCycle} in ~\autoref{ch:loops}
Nested loop thinning \label{Evernote:Max area under histogram}
~\autoref{lst:largest-rectangle-histogram}.


~\autoref{lst:largest-rectangle-histogram} is an example of such usage for solving the
`largest rectangle in histogram' problem
\@\footnote{\url{https://leetcode.com/problems/largest-rectangle-in-histogram}}.

\lstinputlisting[label={lst:largest-rectangle-histogram},
caption={Maximum area under histogram}]
{embedded-loops/histogram.java}
To see the benefit of using the `one-more-cycle' technique, try iterate \code{h.length} times and
do the rest in post-loop logic and testify yourself for the duplicate code between loop
logic and post-loop processing.

\section{Case Study: II}\label{sec:nested-loop-skyline}

%\label{Evernote:Calculate the silhouette or skyline of a group of rectangles on the horizon}
\lstinputlisting[label={lst:silhouette-skyline},
caption={Loop thinning example for skyline calculation.}]
{embedded-loops/silhouette.py}

\label{Evernote:Server peak load problem}
\label{Evernote:for-loop Loop compactification examples}


\section{Case Study: III}\label{sec:nested-loop-quick-sort}

This is a classical example for demonstration of simplicity in programming\cite{Bentley:1986:PP}.
Interested reader may refer to that classical book by Bentley for more detail.

For example, the quick sort algorithm may be implemented as

\begin{clisting}
    \lstinputlisting[label={lst:nestedForLoop},
    caption={Code snippet with nested for loops.}]{embedded-loops/quicksort_outline.java}
\end{clisting}

The for condition and the while condition bear some similarity, so we can combined them into one

Example

\begin{clisting}
    \lstinputlisting[label={lst:nestedForLoop},
    caption={Code snippet with thinned for loops.}]{embedded-loops/quicksort_outline.java}
\end{clisting}

May be made more compact using the above technique as

\begin{clisting}
    \lstinputlisting[label={lst:nestedForLoop},
    caption={Code snippet with thinned for loops.}]{embedded-loops/for_loop_thinned.java}
\end{clisting}


\section{Case Study: IV}\label{sec:nested-loop-numeral-conversion-roman}

\begin{clisting}
    \lstinputlisting[label={lst:nestedForLoop},
    caption={Code snippet with thinned for loops.}]{embedded-loops/int_to_roman_nested_loop.java}
\end{clisting}
\begin{clisting}
    \lstinputlisting[label={lst:nestedForLoop},
    caption={Code snippet with thinned for loops.}]{embedded-loops/int_to_roman_thinned_loop.java}
\end{clisting}


\section{Case Study: V}\label{sec:nested-loop-pattern-matching}

Uncompact:

\begin{clisting}
    \lstinputlisting[label={lst:nestedLoop},
    caption={Code snippet with nested loops.}]{embedded-loops/pattern_match_nested_loop.java}
\end{clisting}

Compact:

\begin{clisting}
    \lstinputlisting[label={lst:thinnedLoop},
    caption={Code snippet with nested loops.}]{embedded-loops/pattern_match_thinned_loop.java}
\end{clisting}


\section{Conclusion}\label{sec:loop-push-up-conclusion}

%What used to be one loop per data structure can actually be achieved with a one-layer loop plus
%cascade conditionals.


Nested loops are hard to understand and shall generally be avoided in production code.
As can be seen, where nested loops occur are mostly related to dynamics adjustment to a
streamlined processing code.
With the `loop thinning' technique discussed in this chapter, many cases of nested loop can be
compressed into single-layer loops.
If this is not possible or the resulting code is harder to understand, one may still consider
refactor inner loops into more meaningful, reusable functions.
